% ===================================================================
%  TABLO No 6 — Evin içi. — Mobilya. — Yemek. — Aydınlatma.
%  Traduction 2026 d'apres texto/io, controlee sur texto/fr et
%  texto/hi. Consigne : voir texto/tr/10-tablo-01.tex.
%
%  Deux notes, toutes deux marquees « (*) », et deux appels : celui de
%  « babo » au bloc c1-12-1, celui de « lambrequino » au bloc c2-01-1.
%  L'ordre les departage, comme dans les autres colonnes.
%
%  LA FEMME DE CHAMBRE N'A PAS DE GROS PLAN, et c'est normal. L'ido du
%  bloc c1-06-1 porte « la chambristino (6) », le francais « la femme
%  de chambre (16) » : c'est l'ido qui se trompe, son propre numero 6
%  etant le savon du bloc c1-02-1. Le turc suit donc le francais et
%  ecrit (16) ; la planche ne porte pas de 16 releve, et le renvoi
%  reste du texte ordinaire. C'est un des trois ecarts declares dans
%  outils/renvoji.py.
% ===================================================================

%%K t06-apar-1 apar
\VUsaut{6.0mm}
\VUtitre{15.0pt}{60}{TABLO No 6}
\VUsaut{1.4mm}
\VUtitre{11.4pt}{0}{Evin içi, mobilya, yemek,}
\VUsaut{0.4mm}
\VUtitre{11.4pt}{0}{aydınlatma.}
\VUsaut{2.2mm}

%%K t06-apar-2 apar
Ev bitmiştir. Jan'ın amcası, düzenini bildiğimiz yeni evine yerleşmiştir.
Şimdi onu nasıl döşediğine bakalım. Önce şunu göreceğiz :

%%K t06-tit-1 sub
\VUpk{10.2pt}{40}{Yatak odası.}
\VUsaut{3.0mm}

%%K t06-c1-01-1 p
1. --- Biraz vakitsiz geldik, çünkü hizmetçi \VUgras{odayı} temizlemekle
uğraşıyor.

%%K t06-c1-02-1 p
2. --- \VUgras{Lavabo masasının} \textsuperscript{(1)} üstünde yıkanmaya,
taranmaya ve hazırlanmaya yarayan şeyler şurada burada durur. Şunlardır :
\VUgras{leğen} \textsuperscript{(2)} ile \VUgras{su ibriği}
\textsuperscript{(3)}, \VUgras{saç fırçası} \textsuperscript{(4)},
\VUgras{tarak} \textsuperscript{(5)}, \VUgras{sabun} \textsuperscript{(6)} ile
\VUgras{sünger} \textsuperscript{(7)}, ve sonunda ağız çalkalama suyunun küçük
bir \VUgras{şişesi} \textsuperscript{(8)}.

%%K t06-c1-03-1 p
3. --- Yerde, masanın önünde, işte kirli suyun \VUgras{kovası}
\textsuperscript{(9)}.

%%K t06-c1-04-1 p
4. --- \VUgras{Dış odada} \textsuperscript{(10)} birkaç \VUgras{askı}
\textsuperscript{(13)} ile iki \VUgras{şemsiye} \textsuperscript{(11)}
görüyoruz, ki \VUgras{şemsiyelikte} \textsuperscript{(12)} durur.

%%K t06-c1-05-1 p
5. --- Kapının öte yanında \VUgras{aynalı dolap} \textsuperscript{(14)}
durur; içinde \VUgras{giysiler} \textsuperscript{(15)} saklanır.

%%K t06-c1-06-1 p
6. --- \VUgras{Hizmetçi} \textsuperscript{(16)} \VUgras{yatağı}
\textsuperscript{(17)} yapıyor; bu vesileyle ayrı ayrı bölümlerini görelim.
\VUgras{Karyolanın} \textsuperscript{(20)} üstünde iyi bir \VUgras{yaylı
şilte} \textsuperscript{(21)} vardır; Jüstin birazdan onun üstüne yün
\VUgras{döşeği} \textsuperscript{(22)} serecek. Sonra sandalyelerden iki
\VUgras{çarşaf} \textsuperscript{(23)} ile \VUgras{battaniyeyi}
\textsuperscript{(24)} alacak. Sonra yatağın başucuna \VUgras{uzun yastığı}
\textsuperscript{(25)} ile \VUgras{yastığı} \textsuperscript{(26)} koyacak, ki
şimdi \VUgras{yorganla} \textsuperscript{(27)} birlikte \VUgras{yatak önü
kilimi} \textsuperscript{(32)} üstünde durur. \VUgras{Perdelere}
\textsuperscript{(19)} ve \VUgras{gölgeliğe} \textsuperscript{(18)} tüylü
süpürgeyi gezdirir, ve iş yarılanır.

%%K t06-c1-07-1 p
7. --- Sonra \VUgras{komodini} \textsuperscript{(28)} biraz derleyip toplaması,
\VUgras{çalar saati} \textsuperscript{(29)} kurması, \VUgras{gece lambasını}
\textsuperscript{(30)} hazırlaması, ve \VUgras{mumu} \textsuperscript{(31)}
alıp şamdanı temizlemesi gerekecek.

%%K t06-tit-2 sub
\VUsaut{5.0mm}
\VUpk{10.2pt}{40}{Dikiş sepeti ve şömine takımı.}
\VUsaut{3.0mm}

%%K t06-c1-08-1 p
8. --- \VUgras{Dikiş sepetinin} \textsuperscript{(34)} de derlenmeye pek
ihtiyacı vardır; \VUgras{sehpanın} \textsuperscript{(33)} yanında durur.
\VUgras{Nakışı} \textsuperscript{(35)} katlaması, \VUgras{dikiş masasına}
\textsuperscript{(38)} \VUgras{yüksüğü}, \VUgras{ipliği}
\textsuperscript{(36)}, \VUgras{iğneleri} \textsuperscript{(37)} ve
\VUgras{makası} \textsuperscript{(39)} koyması, ve masanın kapağını
\VUgras{anahtarla} \textsuperscript{(40)} kilitlemesi gerekecek.

%%K t06-c1-09-1 p
9. --- Ortada duran \VUgras{yuvarlak masanın} \textsuperscript{(41)} üstünde
Jüstin \VUgras{sürahinin} \textsuperscript{(42)} suyunu değiştirecek.
\VUgras{Şekerliğe} \textsuperscript{(43)} \VUgras{şeker} koyacak,
\VUgras{şeker maşasını} \textsuperscript{(44)} temizleyecek ve
\VUgras{elbise fırçasını} \textsuperscript{(46)} kaldıracak.

%%K t06-c1-10-1 p
10. --- Odanın sağ yanı ona daha az iş verecek, çünkü \VUgras{gündüz yatağı}
\textsuperscript{(47)} ile \VUgras{minderi} \textsuperscript{(48)}
yerlerindedir, ve hava soğuk olmadığından \VUgras{şömine}
\textsuperscript{(49)} takımında hiçbir şey oynatılmamıştır. \VUgras{Kürek}
\textsuperscript{(50)}, \VUgras{maşa} \textsuperscript{(51)},
\VUgras{odun ayakları} \textsuperscript{(55)} ile \VUgras{kül perdesi}
\textsuperscript{(56)} temizdir; \VUgras{ateş perdesi} \textsuperscript{(54)}
oynatılmamıştır ve \VUgras{odun sandığı} \textsuperscript{(52)}
\VUgras{çıralarla} \textsuperscript{(53)} doludur.

%%K t06-noto-1 noto
\VUnotes{143.2mm}{%
(*) Babo = küçük çocuk; ingilizce \textit{baby}, fransızca \textit{bébé},
italyanca \textit{bambino}.%
}

%%K t06-noto-2 noto
\VUnotes{143.2mm}{%
(*) Lambrequino = fransızca \textit{lambrequin}, ispanyolca
\textit{lambrequines}, portekizce \textit{lambrequins}, almanca ile
ingilizcede de \textit{lambrequins} --- yani almanca, ingilizce, fransızca,
portekizce ve ispanyolcada, hepsinde bir.%
}

%%K t06-c1-11-1 p
11. --- Daha \VUgras{şömine saatini} \textsuperscript{(57)},
\VUgras{şamdanları} \textsuperscript{(58)} ve \VUgras{süs tabaklarını}
\textsuperscript{(59)} silkelemesi, sonra da \VUgras{aynayı}
\textsuperscript{(60)} silmesi gerekecek.

%%K t06-c1-12-1 p
12. --- Çocuğun \VUgras{beşiğine} \textsuperscript{(61)} gelince (o babo (*)),
o çoktan hazırdır.

%%K t06-tit-3 sub
\VUsaut{5.0mm}
\VUpk{10.2pt}{40}{Salon.}
\VUsaut{3.0mm}

%%K t06-c2-01-1 p
1. --- Şimdi işte salon. Pek güzel bir \VUgras{odadır}. Sağ köşedeki şömine
ince bir \VUgras{heykelcik} \textsuperscript{(1)} ile iki güzel \VUgras{vazoyla}
\textsuperscript{(3)} süslüdür, ve kadifeden bir \VUgras{saçakla}
\textsuperscript{(2)} (*) örtülüdür.

%%K t06-c2-02-1 p
2. --- Şöminenin kıvılcımları halıyı tutuşturmasın diye, önüne bakırdan bir
\VUgras{ateş parmaklığı} \textsuperscript{(4)} konmuştur.

%%K t06-c2-03-1 p
3. --- Yanında şık bir \VUgras{yazı masası} \textsuperscript{(5)} açık durur.
\VUgras{Ev sahibesi} \textsuperscript{(23)} mektuplarını burada yazar.

%%K t06-c2-04-1 p
4. --- Pencerede \VUgras{tül perdeler} \textsuperscript{(6)} vardır, ki
\VUgras{kordonlar} \textsuperscript{(8)} onları bağlı tutar, ve o kordonlar
\VUgras{çengellere} \textsuperscript{(7)} takılıdır; üstlerinde de kalın kumaş
perdeler, ki üstlerinde bir \VUgras{saçak pervazı} \textsuperscript{(9)}
vardır.

%%K t06-c2-05-1 p
5. --- Pencerenin oyuğunda, bir \VUgras{saksılıkta} \textsuperscript{(11)}
yeşil bir \VUgras{bitki} \textsuperscript{(10)} vardır, görkemli bir palmiye,
ki yaprakları \VUgras{gaz lambasına} \textsuperscript{(12)} neredeyse erişir.

%%K t06-c2-06-1 p
6. --- \VUgras{Abajuru} \textsuperscript{(13)} Marya yapmıştır,
o alımlı \VUgras{genç kız} \textsuperscript{(14)}, ki \VUgras{piyanonun}
\textsuperscript{(15)} başında oturur. \VUgras{Taburede} \textsuperscript{(16)} dimdik
oturmuş bir parça çalışıyor. Pek beceriklidir ve pek düzenlidir, ki
\VUgras{nota dolabından} \textsuperscript{(17)} anlaşılır : içinde tam bir
düzen vardır.

%%K t06-tit-4 sub
\VUsaut{5.0mm}
\VUpk{10.2pt}{40}{Yaramaz.}
\VUsaut{3.0mm}

%%K t06-c2-07-1 p
7. --- Erkek kardeşi Pol bir \VUgras{kitap} \textsuperscript{(19)} arıyor,
\VUgras{kitaplıkta} \textsuperscript{(18)}. Yerinde duramayan ve hiç
gemlenmeyen bir \VUgras{oğlandır} \textsuperscript{(20)}. Pek yükseğe konmuş
kitaba erişmek için kitaplığa asılıyor ve üstünde duran \VUgras{heykeli}
\textsuperscript{(21)} düşürme tehlikesine giriyor.

%%K t06-c2-08-1 p
8. --- Bu yaramazlığı yapmak için annesinin bir arkadaşıyla konuşmasından
yararlanıyor; o arkadaş, birkaç gün kendileriyle geçirmeye gelmiş bir
\VUgras{hanımdır} \textsuperscript{(22)}. Anne \VUgras{kanepede}
\textsuperscript{(24)} oturur, \VUgras{koltuğun} \textsuperscript{(25)}
yanında, arkadaşı da bir \VUgras{sandalyede} \textsuperscript{(27)}, ki yalnız
\VUgras{arkalığını} \textsuperscript{(26)} görüyoruz.

%%K t06-c2-09-1 p
9. --- Duvara asılı büyük \VUgras{çerçevede} \textsuperscript{(30)} bir
akrabanın \VUgras{portresi} \textsuperscript{(29)} vardır. Masada duran
\VUgras{albümde} \textsuperscript{(32)} ailenin geri kalanının
\VUgras{fotoğrafları} \textsuperscript{(33)} toplanmıştır. Çocuklar biraz önce
onu karıştırıyordu, çünkü \VUgras{sandalyeler} \textsuperscript{(31)} hâlâ
masanın çevresinde toplu durur.

%%K t06-c2-10-1 p
10. --- Porselenden \VUgras{çin vazosunu} \textsuperscript{(34)}, ki
\VUgras{kâğıt bıçağının} \textsuperscript{(35)} yanında durur, Marya ad
gününde almıştı, ve odanın köşesini dolduran \VUgras{katlanır perdeyi}
\textsuperscript{(36)} amcası Çin'den getirmişti.

%%K t06-c2-11-1 p
11. --- Güzel \VUgras{tablolarla} \textsuperscript{(37)} pırıl pırıl, ki
\VUgras{duvara} \textsuperscript{(42)} asılıdır, tavandaki \VUgras{avizeyle}
\textsuperscript{(38)} aydınlık, ki \VUgras{elektrik lambaları}
\textsuperscript{(39)} akşamları öyle sevinçli parlar, bu salon pek hoş
görünür. Döşemeye serili yumuşak \VUgras{halı} \textsuperscript{(40)} ile öyle
işe yarayan \VUgras{elektrik zili} \textsuperscript{(41)} onu büsbütün rahat
kılar.

%%K t06-tit-5 sub
\VUsaut{3.6mm}
\VUpk{10.2pt}{40}{Yemek odası.}
\VUsaut{3.0mm}

%%K t06-c3-01-1 p
1. --- Yemek çanı çalmıştır. Marya dışında bütün aile sofranın çevresinde
toplanmıştır. Yemek odası iyi bir toprak \VUgras{sobayla} \textsuperscript{(1)} hoş bir
biçimde ısınmıştır, ki \VUgras{borusu} \textsuperscript{(2)} incedir.

%%K t06-c3-02-1 p
2. --- Bu odada eşya çok değildir. Duvar boyunca, \VUgras{sehpanın}
\textsuperscript{(4)} üstünde, süslü küçük bir \VUgras{çeşme}
\textsuperscript{(3)} asılıdır. Yanında \VUgras{büfe} \textsuperscript{(5)}
vardır; üstüne soğuk yemekler, tatlılar ve şu sırada gerekmeyen takımlar
konur.

%%K t06-c3-03-1 p
3. --- Nitekim üst gözde bir \VUgras{kızarmış tavuk} \textsuperscript{(7)}
görüyoruz, bir \VUgras{balık} \textsuperscript{(8)} ve bir \VUgras{kâse}
\textsuperscript{(10)} \VUgras{meyve} \textsuperscript{(9)}. Altta
\VUgras{peynir} \textsuperscript{(11)}, \VUgras{ekmek} \textsuperscript{(12)}
ve \VUgras{yağ ile sirke takımı} \textsuperscript{(13)} vardır; içinde salatayı
donatacak her şey durur : yağ, sirke, \VUgras{tuz} \textsuperscript{(14)} ve
\VUgras{biber} \textsuperscript{(15)}. Sonunda en alt gözde bir
\VUgras{pasta} \textsuperscript{(17)} vardır, \VUgras{hardal kabının}
\textsuperscript{(16)} yanında.

%%K t06-c3-04-1 p
4. --- Yemek odasının saati, ki ona \VUgras{duvar saati} \textsuperscript{(20)}
denir, pek tuhaf biçimlidir. Tahtadan bir çerçeveye oturtulmuştur; o çerçevede
bir \VUgras{barometre} \textsuperscript{(19)} ile bir \VUgras{termometre}
\textsuperscript{(18)} de vardır.

%%K t06-tit-6 sub
\VUsaut{4.6mm}
\VUpk{10.2pt}{40}{Yemeğin verilişi.}
\VUsaut{3.0mm}

%%K t06-c3-05-1 p
5. --- Odanın çevresinde duvarlara cilalı meşeden bir \VUgras{tahta lambri}
\textsuperscript{(21)} çakılmıştır. Kumaştan bir \VUgras{kapı perdesi}
\textsuperscript{(22)} verandaya giden kapıyı iyice örter. Yemekten sonra aile
kahvesini orada içmeye gidecek. \VUgras{Fincanlar} \textsuperscript{(24)} ile
\VUgras{tabaklar} \textsuperscript{(25)} \VUgras{porselen dolabının}
\textsuperscript{(23)} üstüne şimdiden dizilmiştir.

%%K t06-c3-06-1 p
6. --- Sofranın üstünde tavandan bir \VUgras{asma lamba} \textsuperscript{(26)}
sarkar; pek güçlü ışığı akşamları bütün yemek odasını doldurur.

%%K t06-c3-07-1 p
7. --- Şimdi \VUgras{yemek sofrasının} \textsuperscript{(27)} kendisine
bakalım. Aile geldiğinde sofra kurulmuştu. \VUgras{Örtünün}
\textsuperscript{(28)} üstünde her yiyenin yerinde bir \VUgras{tabak}
\textsuperscript{(33)}, bir \VUgras{peçete} \textsuperscript{(29)}, bir
\VUgras{kaşık} \textsuperscript{(34)}, bir \VUgras{çatal}
\textsuperscript{(35)}, bir \VUgras{bıçak} \textsuperscript{(36)} ve bir
\VUgras{bardak} \textsuperscript{(32)} duruyordu. Şarap \VUgras{şişeleri}
\textsuperscript{(30)} ile bir \VUgras{su sürahisi} \textsuperscript{(31)} de
vardı.

%%K t06-c3-08-1 p
8. --- \VUgras{Uşak} \textsuperscript{(39)} önce \VUgras{çorba kâsesini}
\textsuperscript{(37)} getirdi, ki içinde hâlâ biraz çorba tütüyor. Şimdi ilk
\VUgras{yemeği} \textsuperscript{(38)} veriyor, etten olanı. Sonra adlarını
saydığımız yemekleri getirecek; ve sonunda, \VUgras{tatlıdan} sonra,
efendilerin verandaya geçmiş olmasından yararlanıp sofrayı boşaltacak ve yemek
odasını yeniden düzeltecek.

%%K t06-c3-08-2 p
\textit{Not.} --- \textit{Yemeğe bağlı günlük sözler burada ancak kabaca
verilmiştir. Bu liste « otel » (No 13) ile « pazar » (No 15) adlı iki tabloda
tamamlanacaktır.}

%%K t06-tit-7 sub
\VUsaut{4.6mm}
\VUpk{10.2pt}{40}{Mutfak.}
\VUsaut{3.0mm}

%%K t06-c4-01-1 p
1. --- Aralık kapıdan göz attığımız mutfakta görünüşü hoş bir dağınıklık
buluyoruz. \VUgras{Ocağın} \textsuperscript{(1)} önünde, ki orada
\VUgras{ateş} \textsuperscript{(3)} çıtırdıyor, \VUgras{kızartma çevirgeci}
\textsuperscript{(5)} kurulmuştur. İyi bir \VUgras{kızartma}
\textsuperscript{(7)} \VUgras{şişin} \textsuperscript{(6)} üstündedir ve
pişmek üzeredir.

%%K t06-c4-02-1 p
2. --- Az önce kullanılmış olan \VUgras{körük} \textsuperscript{(8)} ile
\VUgras{kömür küreği} \textsuperscript{(9)}, \VUgras{odunlarla}
\textsuperscript{(10)} birlikte yerde bırakılmıştır.

%%K t06-c4-03-1 p
3. --- Şöminenin üst yanı derlidir : \VUgras{fener} \textsuperscript{(11)},
\VUgras{kibritlik} \textsuperscript{(13)}, ki içinde \VUgras{kibritler}
\textsuperscript{(12)} vardır, \VUgras{şamdan} \textsuperscript{(14)} ve
\VUgras{el şamdanı} \textsuperscript{(15)}, \VUgras{mumuyla}
\textsuperscript{(16)} birlikte --- hepsi yerindedir. Ama büyük
\VUgras{kömür sepeti} \textsuperscript{(18)}, \VUgras{kömürle}
\textsuperscript{(17)} dolu, ki \VUgras{aşçı kadın} \textsuperscript{(23)} az
önce bodrumdan getirmiştir, mutfağın tam ortasında kalıvermiştir.

%%K t06-c4-04-1 p
4. --- Doğrusu, öğle yemeğini vaktinde hazırlayacaksa zavallının acele etmesi
gerek. Şu sırada \VUgras{pişirme ocağının} \textsuperscript{(19)} üstünde bir
\VUgras{sos} \textsuperscript{(24)} karıştırıyor, yanmasın diye.

%%K t06-c4-05-1 p
5. --- \VUgras{Fırında} \textsuperscript{(20)} bir pasta pişiyor, ki çocukların
akşam çayı için yapmıştır. Ocağın üstünde \VUgras{kepçe} \textsuperscript{(21)}
ile \VUgras{kevgir} \textsuperscript{(22)} asılıdır.

%%K t06-tit-8 sub
\VUsaut{4.0mm}
\VUpk{10.2pt}{40}{Mutfak takımları.}
\VUsaut{3.0mm}

%%K t06-c4-06-1 p
6. --- Odanın arkasında asılı \VUgras{kap kacak takımı} \textsuperscript{(25)}
pek temizdir. Bakırdan güzel \VUgras{tencereler} \textsuperscript{(26)},
\VUgras{süt güğümü} \textsuperscript{(27)}, \VUgras{süzgeç}
\textsuperscript{(28)} ve \VUgras{rende} \textsuperscript{(29)} özenle
ovulmuştur.

%%K t06-c4-07-1 p
7. --- \VUgras{Tuzluk} \textsuperscript{(30)} küçük dolabın üstünde
\VUgras{çaydanlık} \textsuperscript{(31)} ile \VUgras{büyük yağ şişesinin}
\textsuperscript{(32)} yanında durur. \VUgras{Çekmecede}
\textsuperscript{(33)} herhâlde kaşık çatal ile mutfak bıçakları durur.

%%K t06-c4-08-1 p
8. --- \VUgras{Süpürge} \textsuperscript{(34)} ile \VUgras{tüylü süpürge}
\textsuperscript{(35)} bir köşededir. Aşçı kadın az önce \VUgras{kütüğü}
\textsuperscript{(36)} kullanmıştır; onu pencerenin yanında bırakmıştır.

%%K t06-c4-09-1 p
9. --- Bir \VUgras{el havlusu} \textsuperscript{(37)} duvarda asılıdır,
\VUgras{huninin} \textsuperscript{(38)} yanında. Mutfağın bu yanında
\VUgras{şiş ızgarası} \textsuperscript{(39)}, \VUgras{satır}
\textsuperscript{(40)} ve \VUgras{kesme tahtası} \textsuperscript{(41)} durur.
Sonra satırın üstünde, bir \VUgras{rafta} \textsuperscript{(42)},
\VUgras{çömlekler} \textsuperscript{(43)}, \VUgras{güveç}
\textsuperscript{(44)} \VUgras{kapağıyla} \textsuperscript{(45)} birlikte, ve
\VUgras{kazan} \textsuperscript{(46)} vardır. \VUgras{Kızartma tavası}
\textsuperscript{(47)} yanı başında asılıdır.

%%K t06-c4-10-1 p
10. --- Bu köşede yalnız bakırdan \VUgras{musluk} \textsuperscript{(48)}
parıldar. \VUgras{Eviye} \textsuperscript{(49)}, ki üstünde büyük
\VUgras{testi} \textsuperscript{(50)} durur, küçük dolabıyla birlikte pek işe
yarayacaktır; orada \VUgras{sirke fıçısı} \textsuperscript{(51)},
\VUgras{musluğuyla} \textsuperscript{(52)} birlikte, \VUgras{çöp sandığı}
\textsuperscript{(53)} ve \VUgras{kirli kaplar} \textsuperscript{(54)} durur.

%%K t06-tit-9 sub
\VUsaut{4.0mm}
\VUpk{10.2pt}{40}{Mutfakta duran başka şeyler.}
\VUsaut{3.0mm}

%%K t06-c4-11-1 p
11. --- Şu büyük \VUgras{tekne} \textsuperscript{(55)}, ki içinde
\VUgras{sebzeler} \textsuperscript{(58)} yıkanır, ve dökme demirden
\VUgras{kazan} \textsuperscript{(56)}, ki \VUgras{tuğla döşemenin}
\textsuperscript{(57)} üstünde durur, herhâlde o dolabındır.

%%K t06-c4-12-1 p
12. --- Aşçı kadının ocağı eksik değildir, çünkü işte masanın köşesinde bir
\VUgras{gaz ocağı} \textsuperscript{(59)}, üstünde küçük bir tencerede su
ısınıyor. Ondan çıkan \VUgras{buhar} \textsuperscript{(60)} kaynamakta
olduğunu gösterir. Bu su herhâlde kahve içindir, çünkü masanın öteki ucunda
işte \VUgras{kahve kabı} \textsuperscript{(64)} ile \VUgras{kahve değirmeni}
\textsuperscript{(63)}.

%%K t06-c4-13-1 p
13. --- Gaz ocağı \VUgras{gaz musluğuna} \textsuperscript{(62)} bağlıdır, uzun
bir \VUgras{lastik boru} \textsuperscript{(61)} aracılığıyla.

%%K t06-c4-14-1 p
14. --- Aşçı kadına yardıma gelen \VUgras{gündelikçi} \textsuperscript{(66)},
akşam yemeğinin sebzelerini hazırlıyor. Soyup yıkadıktan sonra, pişirme vakti
gelene dek onları \VUgras{kaba} \textsuperscript{(65)} koyacak.

%%K t06-tit-10 sub
\VUsaut{4.0mm}
{\centering\fontsize{10.2pt}{10.2pt}\selectfont\textls[40]{\scshape BANYO.} \textit{(Yıkanma odası.)}\par}
\VUsaut{3.0mm}

%%K t06-c5-01-1 p
1. --- Evin başlıca odalarını öğrenmek için geriye bir tek göz atmak kalıyor :
banyonun döşenişini görmek. Uzun sürmeyecek, çünkü büyük değildir, ve hemen
göreceğiz ki \VUgras{banyo küvetinin} \textsuperscript{(1)} yanında iyi bir
\VUgras{su ısıtıcısı} \textsuperscript{(2)} vardır, \VUgras{sabunluk}
\textsuperscript{(3)} yıkananın elinin erişeceği yerdedir, ve
\VUgras{havlulukla} \textsuperscript{(5)} \VUgras{havluları}
\textsuperscript{(4)} kurutmak kolay oluyordur.

%%K t06-c5-02-1 p
2. --- Bu odanın duvarları parlak \VUgras{fayanslarla} \textsuperscript{(6)}
kaplıdır; bu da bir \VUgras{duş} \textsuperscript{(7)} koymayı olanaklı kıldı,
kullanılırken sıçrayan suların onları bozmasından korkmadan. Ayak yıkamaya
yarayan çinkodan küçük bir \VUgras{leğen} \textsuperscript{(8)} döşenişi
tamamlar.
\VUsaut{6.0mm}
\VUfilet{22mm}
